\chapter{فضاءات هيلبرت}\label{ch:b3:hilbert}

\omterm{def:b3:hilbert:inner}{فضاء هيلبرت} فضاءُ باناخ معيارُه آتٍ من \omterm{def:b3:hilbert:inner}{جداء سلّمي} --- وهذا
البناء الإضافي الوحيد يعيد، في البُعد اللانهائي، جُلّ الهندسة
الإقليدية: فالإسقاطات المتعامدة موجودة، وكل مؤثّرة خطية متصلة
هي \omterm{def:b3:hilbert:inner}{جداء سلّمي} بمتجهة ثابتة (ريس)، والأسس المتعامدة المتجانسة
تنشر كل متجهة في متسلسلة متقاربة بمحاسبة فيثاغورية (\omterm{thm:b3:hilbert:parseval}{بارسيفال}).
وذروة الفصل وفاءٌ بدَين: أن النظام المثلثي أساسٌ متعامد متجانس
لفضاء $L^2$، ومنه تصح متطابقة \omterm{thm:b3:hilbert:parseval}{بارسيفال} من أجل \emph{كل} \omterm{def:b3:lebesgue:l1}{دالة
قابلة للمكاملة} بالمربّع --- وهي العبارة التي لم تستطع السنة
الجامعية 2 أن تبرهن عليها إلا من أجل الدوال من الصنف
$\mathcal C^1$ بالقطع. وننهي بلاكس--ميلغرام، المبرهنة المساعدة
التي هي حصان المقاربة التغيّرية للمعادلات التفاضلية.

وفي كل ما يلي، يكون $H$ فضاءً متجهيًّا على $K = \R$ أو $\C$.

\section{الجداءات السلّمية؛ مبرهنة الإسقاط}

\begin{definition}\label{def:b3:hilbert:inner}
\emph{الجداء السلّمي}\index{جداء سلّمي} هو تطبيق $\langle
\cdot,\cdot\rangle \colon H\times H \to K$، خطي في
المتغيّر الثاني، يحقق $\langle y, x\rangle =
\overline{\langle x, y\rangle}$ و $\langle x, x\rangle > 0$ من أجل $x \neq 0$. وهو
يستحث المعيار $\norm x = \langle x,
x\rangle^{1/2}$، و\emph{متراجحة كوشي--شوارتز} $\abs{\langle x, y\rangle} \leq \norm x\norm y$ (وبرهانها
--- بالمميِّز --- في السنة الجامعية 2 لم يتغيّر)،
و\emph{قاعدة متوازي الأضلاع}
\[
\norm{x + y}^2 + \norm{x - y}^2 = 2\norm x^2 + 2\norm y^2 .
\]
و\emph{فضاء هيلبرت}\index{فضاء هيلبرت} هو فضاءٌ ذو جداء سلّمي
تامٌّ بالنسبة إلى هذا المعيار. أمثلة: $\ell^2$ (\cref{pb:b3:banach:1})،
والمثال الأساسي $L^2(\mu)$ مع $\langle f, g\rangle = \int\bar fg\,\dd\mu$ --- وهو تام بريس--فيشر
(\cref{thm:b3:lp:complete})؛ والجداء السلّمي منتهٍ بكوشي--شوارتز (أي هولدر
عند $p = q = 2$).
\end{definition}

\begin{theorem}[الإسقاط على مجموعة محدّبة مغلقة]
\label{thm:b3:hilbert:projection}
لتكن $C \neq \varnothing$ جزءًا \emph{محدّبًا} مغلقًا من \omterm{def:b3:hilbert:inner}{فضاء هيلبرت} $H$
وليكن $x \in H$. توجد $p_C(x)
\in C$ وحيدة تحقق\index{مبرهنة الإسقاط}
\[
\norm{x - p_C(x)} = d(x, C),
\]
ويميّزها الشرط: $\operatorname{Re}\langle x - p_C(x),\ c -
p_C(x)\rangle \leq 0$ من أجل كل $c \in C$. والتطبيق $p_C$
ليبشيتزي بالثابت $1$.
\end{theorem}

\begin{proof}
ليكن $d = d(x, C)$ ولتكن $(c_n) \subseteq C$ تحقق $\norm{x - c_n}
\to d$. وبقاعدة متوازي
الأضلاع على $x - c_n$ و $x - c_m$:
\[
\norm{c_n - c_m}^2 = 2\norm{x - c_n}^2 + 2\norm{x - c_m}^2 -
4\,\bigl\|x - \tfrac{c_n + c_m}2\bigr\|^2
\leq 2\norm{x{-}c_n}^2 + 2\norm{x{-}c_m}^2 - 4d^2
\]
(إذ يضع التحدّب المنتصف في $C$): فيؤول الطرف الأيمن إلى $0$،
ومنه تكون $(c_n)$ كوشية، وتبلغ نهايتها $p \in C$ (بالإغلاق)
المسافةَ $d$. وأما الوحدانية: فمصغِّران يعطيان، بالمتطابقة
نفسها، $\norm{p - p'}^2 \leq 2d^2 + 2d^2 - 4d^2 = 0$.

وأما التمييز: فمن أجل $c \in C$ و $t \in \intoc01$، تكون
المتجهة $p + t(c - p) \in C$، ومنه
\[
d^2 \leq \norm{x - p - t(c-p)}^2
= d^2 - 2t\operatorname{Re}\langle x - p, c - p\rangle +
t^2\norm{c-p}^2 ;
\]
ثم نقسم ونجعل $t \to 0^+$: $\operatorname{Re}\langle x - p, c -
p\rangle \leq 0$. وبالعكس تعطي هذه المتراجحة
أن $\norm{x -
c}^2 = \norm{x - p}^2 - 2\operatorname{Re}\langle x - p, c -
p\rangle + \norm{p - c}^2 \geq \norm{x-p}^2$. وأما الليبشيتزية: فمن أجل $x, y$ ذواتَي الإسقاطين
$p, q$، نجمع المتراجحتين التغيّريتين (مع $c = q$ و $c = p$ على
التوالي): $\operatorname{Re}\langle x - y - (p - q), p - q\rangle \geq
0$، ومنه $\norm{p - q}^2 \leq \operatorname{Re}\langle x - y, p -
q\rangle \leq \norm{x - y}\norm{p - q}$.
\end{proof}

\begin{theorem}[التفكيك المتعامد]
\label{thm:b3:hilbert:decomposition}
ليكن $F$ \emph{فضاءً جزئيًّا مغلقًا} من $H$. عندئذٍ يكون $p_F$
خطيًّا، و $x - p_F(x) \perp F$ من أجل كل $x$، و\index{متمّم متعامد}
\[
H = F \oplus F^\perp,
\qquad F^\perp = \{y : \langle y, f\rangle = 0\ \forall f\in
F\},
\qquad (F^\perp)^\perp = F .
\]
وأما من أجل فضاء جزئي عام، فإن $(F^\perp)^\perp = \bar F$؛ وعلى وجه الخصوص يكون $F$
كثيفًا إذا وفقط إذا كان $F^\perp = \{0\}$.
\end{theorem}

\begin{proof}
من أجل فضاء جزئي، يفرض التمييز التغيّري مع $c =
p_F(x) \pm f$ (حيث
$f \in F$، بالإشارتين، ومع $\iu f$ في الحالة العقدية) أن
$\langle x - p_F(x), f\rangle = 0$: أي إن الباقي متعامد مع $F$. والتفكيك $x = p_F(x) + (x
- p_F(x))$ مع $F \cap F^\perp = \{0\}$
($\langle y, y
\rangle = 0$)؛ وتنتج خطية $p_F$ من وحدانية هذه التفكيكات (إذ الطرفان
خطيان فيها). ويصح $(F^\perp)
^\perp \supseteq F$ دائمًا؛ وبالعكس إذا كان $x \perp F^\perp$، كتبنا
$x = f + g$: فيكون $g = x - f \in F^\perp$ و $\langle g,
g\rangle = \langle x, g\rangle - \langle f, g\rangle = 0$: ومنه $x
= f \in F$. وأما من أجل
فضاء جزئي عام $F$: فإن $F^\perp = \bar
F^{\,\perp}$ (باتصال \omterm{def:b3:hilbert:inner}{الجداء السلّمي})، ومنه
$(F^\perp)^\perp = \bar F$ بالحالة المغلقة؛ والكثافة تكافئ $\bar F = H$ التي تكافئ
$F^\perp = 0$.
\end{proof}

\begin{theorem}[تمثيل ريس]\label{thm:b3:hilbert:riesz}
من أجل كل مؤثّرة خطية متصلة $\varphi \in H'$ يوجد $a \in H$
وحيد يحقق\index{مبرهنة ريس في التمثيل}
\[
\varphi(x) = \langle a, x\rangle \quad (x \in H),
\qquad \norm\varphi_{H'} = \norm a .
\]
\end{theorem}

\begin{proof}
إذا كان $\varphi = 0$: فإن $a = 0$. وإلا فإن $F = \ker\varphi$ فضاء جزئي
مغلق فعلي؛ نختار $u \in F^\perp$ يحقق $\norm u = 1$
(\cref{thm:b3:hilbert:decomposition}: إذ $F^\perp \neq 0$ لأن $F \neq H$). ومن أجل أي
$x$، تكون المتجهة $\varphi(x)u -
\varphi(u)x \in \ker\varphi$، ومنه فهي $\perp u$:
\[
0 = \langle u, \varphi(x)u - \varphi(u)x\rangle
= \varphi(x) - \varphi(u)\langle u, x\rangle :
\qquad \varphi(x) = \langle
\overline{\varphi(u)}\,u,\ x\rangle .
\]
ومنه تصلح $a = \overline{\varphi(u)}u$. وأما الوحدانية: فمن $\langle a
- a', x\rangle = 0$ من أجل كل $x$،
باختبار $x = a - a'$. وأما النظائم: فإن $\abs{\varphi(x)} \leq \norm a\norm x$ (بكوشي--شوارتز)
مع التساوي عند $x = a$.
\end{proof}

\begin{example}[إسقاط، محسوبٌ حتى النهاية]
\label{ex:b3:hilbert:projectionexample}
في $H = L^2(\intcc01)$، ما أفضل تقريب للدالة $f(x) = x^2$ بدالة تآلفية؟
الفضاء الجزئي $F =
\operatorname{Vect}(1, x)$ مغلق (لأنه منتهي البُعد)، وتميّز $p_F(f) = a + bx$
بتعامد الباقي مع $1$ ومع $x$:
\[
\int_0^1(x^2 - a - bx)\,\dd x = 0,
\qquad
\int_0^1x\,(x^2 - a - bx)\,\dd x = 0,
\]
أي $\frac13 = a + \frac b2$ و $\frac14 = \frac a2 +
\frac b3$: ومنه $a = -\frac16$ و $b = 1$. فيكون
$p_F(x^2) = x -
\frac16$، ويكون الخطأ
\[
d(f, F)^2 = \int_0^1\Bigl(x^2 - x + \frac16\Bigr)^2\dd x =
\frac1{180},
\qquad d(f, F) = \frac1{6\sqrt5} .
\]
وملاحظتان جديرتان بالاستيعاب. أولًا، ليس الحساب سوى جملة خطية
من الرتبة $2\times2$ --- أي \emph{المعادلات الناظمية}؛ ومن أجل
أساس الوحيدات الحدّية تكون مصفوفتها $\bigl(\frac1{i+j+1}\bigr)$ هي مصفوفة هيلبرت
السيّئة التكييف بشهرتها، والعلاج هو التعامد أولًا (كثيرات حدود
لوجاندر، \cref{pb:b3:hilbert:1}). وثانيًا، أفضل تقريب \emph{منتظم} للمقدار
$x^2$ بالدوال التآلفية مختلفٌ ($x - \frac18$، بالتذبذب
المتساوي): فلكل معيار هندسته، ولا يجيب بجملة خطية إلا المعيار
الهيلبرتي.
\end{example}

\section{الأسس المتعامدة المتجانسة}

\begin{definition}\label{def:b3:hilbert:onb}
تكون العائلة $(e_i)_{i\in I}$ \emph{متعامدة متجانسة} إذا كان
$\langle
e_i, e_j\rangle = \delta_{ij}$، وتكون \emph{أساسًا هيلبرتيًّا}\index{أساس هيلبرتي} (أي
أساسًا متعامدًا متجانسًا) إذا كانت فوق ذلك تركيباتها الخطية
المنتهية كثيفةً في $H$ (أي إن العائلة \emph{شاملة}). ونعالج
الحالة القابلة للعدّ $I = \N$، وهي تغطي بغرام--شميدت كل فضاء $H$
\emph{قابل للفصل} (\cref{prop:b3:hilbert:gramschmidt}).
\end{definition}

\begin{theorem}[بسل، بارسيفال]\label{thm:b3:hilbert:parseval}
لتكن $(e_n)_{n\in\N}$ متعامدة متجانسة في $H$، ولتكن $c_n(x) =
\langle e_n, x\rangle$.
\begin{enumerate}
  \item (بسل) $\sum_n\abs{c_n(x)}^2 \leq \norm x^2$، وتتقارب المتسلسلة $\sum_nc_n(x)e_n$ في $H$،
        ومجموعها $p_F(x)$، حيث $F = \overline{\operatorname{Vect}}(e_n)$.
  \item وتتكافأ العبارات التالية: (أ) إن $(e_n)$ \omterm{def:b3:hilbert:onb}{أساس
        هيلبرتي}؛ (ب) $x = \sum_nc_n(x)e_n$ من أجل كل $x$؛ (ج)
        \emph{بارسيفال}\index{متطابقة بارسيفال}: $\norm x^2
        = \sum_n\abs{c_n(x)}^2$ من أجل
        كل $x$؛ (د) المتجهة الوحيدة المتعامدة مع جميع $e_n$
        هي $0$.
  \item وإذا كانت $(e_n)$ \omterm{def:b3:hilbert:onb}{أساسًا هيلبرتيًّا}، فإن $x \mapsto (c_n(x))_n$ تماثلٌ
        متقايس $H \to \ell^2$ (أي إن \emph{كل} \omterm{def:b3:hilbert:inner}{فضاء هيلبرت}
        لانهائي البُعد قابل للفصل «هو» $\ell^2$)، و $\langle x, y\rangle =
        \sum_n\overline{c_n(x)}c_n(y)$.
\end{enumerate}
\end{theorem}

\begin{proof}
(1) من أجل $N$ منتهٍ: $x - \sum_{n\leq N}c_ne_n \perp e_k$ ($k
\leq N$)، ومنه تعطي فيثاغورس أن
$\norm x^2 = \sum_{n\leq
N}\abs{c_n}^2 + \norm{x - \sum_{n\leq N}c_ne_n}^2$: وهي بسل. والمجاميع الجزئية $S_N = \sum_{n\leq N}c_ne_n$ كوشية: $\norm{S_N - S_M}^2 = \sum_{M<n\leq N}\abs{c_n}^2$، وهو
ذيل متسلسلة متقاربة؛ وتقع النهاية في $F$، وتكون $x - \lim S_N
\perp$ كل $e_k$
(بالاتصال)، ومنه فهي $\perp F$: فبوحدانية التفكيك المتعامد،
$\lim S_N = p_F(x)$.

(2) (أ)$\Rightarrow$(ب): لأن $F = H$، ومنه $p_F = \mathrm{id}$.
(ب)$\Rightarrow$(ج): بفيثاغورس عند النهاية ($\norm{S_N}^2
= \sum_{n \leq N}\abs{c_n}^2 \to \norm x^2$).
(ج)$\Rightarrow$(د): فمن $x \perp$ جميع $e_n$ نجد $\norm x^2 =
0$.
(د)$\Rightarrow$(أ): $F^\perp = \{0\}$ (إذ التعامد مع جميع $e_n$ هو
التعامد مع $F$)، ومنه يكون $F$ كثيفًا حسب \cref{thm:b3:hilbert:decomposition}؛ لكن $F$
غلقٌ، فهو مغلق أصلًا: ومنه $F = H$.

(3) التطبيق خطي، ومتقايس حسب (ج) (ومنه فهو متباين)، وغامر:
إذ بإعطاء $(c_n) \in \ell^2$ تتقارب المتسلسلة $\sum
c_ne_n$ (وهي كوشية كما في (1))
إلى صورة عكسية. وأما صيغة \omterm{def:b3:hilbert:inner}{الجداء السلّمي} فهي الاستقطاب انطلاقًا
من (ج)، أو حسابٌ مباشر بالنهاية.
\end{proof}

\begin{proposition}[غرام--شميدت]\label{prop:b3:hilbert:gramschmidt}
لتكن $(x_n)$ متتالية مستقلة خطيًّا. فبالوضع بالتراجع $\tilde e_n = x_n - \sum_{k<n}\langle e_k,
x_n\rangle e_k$
و $e_n = \tilde e_n/\norm{\tilde e_n}$ نحصل على عائلة متعامدة متجانسة $(e_n)$ لها الفضاءات
المولَّدة المنتهية نفسها: $\operatorname{Vect}(e_1, \dots, e_n) = \operatorname{Vect}
(x_1, \dots, x_n)$. ومن ثَمّ فلكل \omterm{def:b3:hilbert:inner}{فضاء هيلبرت}
قابل للفصل (أي ذي جزء كثيف قابل للعدّ) أساسٌ
هيلبرتي.\index{تعامد غرام--شميدت}
\end{proposition}

\begin{proof}
بالتراجع: $\tilde e_n \perp e_k$ (حيث $k < n$) بحكم البناء، و $\tilde e_n \neq 0$ بالاستقلال؛
وتتوافق الفضاءات المولَّدة في كل مرحلة (بتغيير أساس مثلثي). وأما
من أجل $H$ قابل للفصل: فمن متتالية كثيفة نستخرج عائلةً جزئية
مستقلة خطيًّا ذات فضاء مولَّد كثيف (بطرح كل متجهة تقع في الفضاء
المولَّد بسابقاتها --- والفضاء المولَّد لا يتغيّر)، ثم نعامدها:
فتكون النتيجة شاملة.
\end{proof}

\begin{theorem}[النظام المثلثي؛ بارسيفال أخيرًا]
\label{thm:b3:hilbert:fourier}
في $L^2(\intcc{-\pi}\pi)$ مع $\langle f, g\rangle =
\frac1{2\pi}\int_{-\pi}^\pi \bar fg$، تكون العائلة $e_n(t) =
\eu^{\iu nt}$، حيث $n \in \Z$،
\omterm{def:b3:hilbert:onb}{أساسًا هيلبرتيًّا}. ومن ثَمّ، من أجل \emph{كل} $f \in L^2$ ---
وعلى وجه الخصوص كل دالة متصلة بالقطع دورية بالدور $2\pi$ ---
ومع $c_n(f) =
\frac1{2\pi}\int_{-\pi}^{\pi}f(t)\eu^{-\iu nt}\dd t$:
\[
f = \sum_{n\in\Z}c_n(f)\,\eu^{\iu nt} \ \ \text{في} L^2,
\qquad
\frac1{2\pi}\int_{-\pi}^{\pi}\abs f^2 =
\sum_{n\in\Z}\abs{c_n(f)}^2 .
\]
وهذا يبرهن، في عموم تام، على متطابقة \omterm{thm:b3:hilbert:parseval}{بارسيفال} التي سلّمت بها
السنة الجامعية 2.\index{متطابقة بارسيفال (العامة)}
\end{theorem}

\begin{proof}
التعامد والتجانس حسابٌ مباشر (السنة الجامعية 2). وأما
الشمول: فلتكن $f \in L^2$ تحقق $\perp$ جميع $e_n$، أي أن تنعدم
جميع معاملات فورييه لها. والدوال المتصلة الدورية بالدور $2\pi$
كثيفة في $L^2(\intcc{-\pi}\pi)$: إذ $\mathcal
C_c(\intoo{-\pi}\pi)$ كثيفة (\cref{thm:b3:lp:density}(2)) وتمتدّ هذه
الدوال دوريًّا واتصاليًّا. وكثيرات الحدود المثلثية كثيفة
بالمعنى $\norm\cdot_\infty$ بين الدوال المتصلة الدورية (بستون--فايرشتراس،
\cref{cor:b3:complete:weierstrass}(c))، و $\norm\cdot_2 \leq \norm\cdot_\infty$: ومنه فكثيرات الحدود المثلثية كثيفة
في $L^2$. لكن $f \perp$ كل كثير حدود مثلثي، ومنه $f \perp$ فضاء
جزئي كثيف: أي $f
\in (\text{كثيف})^\perp = \{0\}$ (\cref{thm:b3:hilbert:decomposition}). ويخلص
المحك (د) من \cref{thm:b3:hilbert:parseval} إلى النتيجة؛ ويُفكّ المحكان (ب)
و (ج) إلى العبارة المعروضة (بإعادة ترقيم المجموعة القابلة للعدّ
$\Z$؛ والمتسلسلة ذات الطرفين تتقارب تقاربًا غير مشروط --- إذ
تتقارب المجاميع الجزئية على أي عائلة مستنفدة، بحجة الذيل في
$\ell^2$).
\end{proof}

\begin{theorem}[لاكس--ميلغرام]\label{thm:b3:hilbert:laxmilgram}
ليكن $H$ \omterm{def:b3:hilbert:inner}{فضاء هيلبرت} حقيقيًّا ولتكن $a \colon H\times H \to
\R$ ثنائية الخطية
و\emph{متصلة} ($\abs{a(u,v)} \leq M\norm
u\norm v$) و\emph{قسرية} ($a(u, u) \geq \alpha\norm u^2$، حيث
$\alpha > 0$). عندئذٍ يوجد من أجل كل $\varphi \in H'$ عنصرٌ
وحيد $u \in H$ يحقق\index{مبرهنة لاكس--ميلغرام}
\[
a(u, v) = \varphi(v) \qquad \text{من أجل كل} v \in H .
\]
\end{theorem}

\begin{proof}
من أجل $u$ مثبَّتة، تكون $v \mapsto a(u, v)$ خطيةً متصلة: فيعطي ريس عنصرًا
وحيدًا $Au \in H$ يحقق $a(u,v) = \langle Au,
v\rangle$؛ و $A$ خطي مع $\norm{Au} \leq M\norm u$ (بوحدانية
الممثِّلات، ثم بالحد). وأما القسرية: فإن $\alpha\norm u^2 \leq a(u,u) = \langle Au, u\rangle \leq
\norm{Au}\norm u$، ومنه $\norm{Au} \geq \alpha\norm u$:
أي إن $A$ متباين ومجاله مغلق (إذ متتالية صور كوشية $Au_n$ تفرض
أن تكون $u_n$ كوشية). والمجال كثيف: إذ يعطي $w \perp
\operatorname{im}A$ أن $0 = \langle Aw, w\rangle \geq
\alpha\norm w^2$.
ومغلق وكثيف: ومنه فإن $A$ تقابل. وبإعطاء $\varphi$، ليكن $f$
ممثِّله (بريس) وليكن $u = A^{-1}f$: فيكون $a(u, v) = \langle f, v\rangle = \varphi(v)$، وبكيفية
وحيدة ($a(u - u', \cdot) = 0$ والقسرية).
\end{proof}

\begin{remark}\label{rem:b3:hilbert:variational}
حين تكون $a$ متناظرة، يكون حل لاكس--ميلغرام هو المصغِّر الوحيد
لمقدار \emph{الطاقة} $J(v) = \frac12a(v,v) -
\varphi(v)$ (\cref{exo:b3:hilbert:9}): أي وجود حلول للمسائل
التغيّرية بضربة واحدة. وبتطبيق ذلك على فضاءات دوال ملائمة (أي
فضاءات سوبوليف في مقرر لاحق)، تُحلّ مسائل الشروط الحدّية
للمعادلات التفاضلية --- وهو المدخل الحديث إلى المعادلات
التفاضلية الجزئية.
\end{remark}

\section{تمارين}

\begin{exercise}[$\star$]\label{exo:b3:hilbert:1}
(a) برهن على متطابقات الاستقطاب (في الحالة الحقيقية: $4\langle x,
y\rangle = \norm{x+y}^2 - \norm{x-y}^2$؛
وفي العقدية: الصيغة ذات الحدود الأربعة).
(b) برهن على أن $\norm\cdot_1$ على $L^1(\intcc01)$ و $\norm\cdot_\infty$
على $\mathcal C(\intcc01)$ تخرقان قاعدة متوازي الأضلاع: فهذان المعياران لا
يأتيان من أي \omterm{def:b3:hilbert:inner}{جداء سلّمي}.
\end{exercise}

\begin{exercise}[$\star$]\label{exo:b3:hilbert:2}
في $H = L^2(\intcc01)$ (الحقيقي):
(a) احسب إسقاط $f$ على فضاء الدوال الثابتة، وفسّر؛
(b) احسب الإسقاط على $\{g : g = 0 \text{ في كل مكان تقريبًا على}
\intcc0{1/2}\}$؛
(c) احسب $d\bigl(x \mapsto x,\ \operatorname{Vect}(\mathbf
1)\bigr)$.
\end{exercise}

\begin{exercise}[$\star\star$]\label{exo:b3:hilbert:3}
(a) برهن على أنه من أجل فضاء جزئي $F$: يكون $F$ كثيفًا
$\iff$ $F^\perp =
\{0\}$، وأعطِ مثالًا في $\ell^2$ على فضاء جزئي
\emph{فعلي} كثيف (بحيث $F^\perp = 0$ دون $F = H$: فمبرهنة
التفكيك تحتاج فعلًا إلى $F$ مغلقًا).
(b) برهن على أنه إذا كان $x_n \to x$ و $y_n \to y$ بالمعيار،
فإن $\langle x_n, y_n\rangle \to \langle x, y\rangle$، وعيّن موضعين استعمل فيهما الفصل هذا الاتصال.
\end{exercise}

\begin{exercise}[$\star\star$]\label{exo:b3:hilbert:4}
طبّق غرام--شميدت على $1, x, x^2$ في $L^2(\intcc{-1}1)$ (\omterm{def:b3:measure:lebesgueouter}{بقياس لوبيغ}):
لتحصل على \emph{كثيرات حدود لوجاندر} المعيَّرة الثلاثة الأولى،
وتحقق من أنها توافق $\sqrt{n + \frac12}\,P_n$ من أجل كثيرات حدود رودريغ $P_n$ في
\cref{pb:b3:hilbert:1}.
\end{exercise}

\begin{exercise}[$\star\star$]\label{exo:b3:hilbert:5}
طبّق \omterm{thm:b3:hilbert:parseval}{بارسيفال} (\cref{thm:b3:hilbert:fourier}) على $f(t) = t$ و $f(t) = t^2$ على
$\intcc{-\pi}\pi$ --- وهو الآن مشروع من أجلهما (فهما متصلتان، غير أن
المتطابقة كانت تتطلب سابقًا عنايةً من الصنف $\mathcal C^1$
بالقطع عند انقطاع الالتفاف الدوري): لتستعيد
\[
\sum_{n\geq1}\frac1{n^2} = \frac{\pi^2}6,
\qquad
\sum_{n\geq1}\frac1{n^4} = \frac{\pi^4}{90} .
\]
\end{exercise}

\begin{exercise}[$\star\star$]\label{exo:b3:hilbert:6}
(a) جد $a \in L^2(\intcc01)$ تحقق $\int_0^{1/2}f =
\langle a, f\rangle$ من أجل كل $f$؛ واحسب
$\norm\varphi$ لهذه المؤثّرة.
(b) برهن على أن التقييم $f \mapsto f(\frac12)$، المعرَّف على الفضاء الجزئي
$\mathcal C(\intcc01) \subseteq
L^2(\intcc01)$، \emph{غير} متصل بالنسبة إلى $\norm\cdot_2$: فلا يوجد
أي ممثِّل بمعنى ريس (إذ التقييم ليس مفهومًا من مفاهيم $L^2$).
\end{exercise}

\begin{exercise}[$\star\star\star$]\label{exo:b3:hilbert:7}
ليكن $H$ قابلًا للفصل ذا \omterm{def:b3:hilbert:onb}{أساس هيلبرتي} $(e_n)$، ولتكن $(x_k)$
متتالية محدودة.
(a) برهن على أن متتالية جزئية ما تتقارب \emph{تقاربًا ضعيفًا}:
أي يوجد $x$ يحقق $\langle y, x_{k_j}\rangle \to \langle y,
x\rangle$ من أجل كل $y \in H$. \emph{(بالاستخراج
القطري على المعاملات $\langle e_n, x_k\rangle$؛ ثم ركّب $x$ بواسطة بسل والحد
المنتظم للنظائم.)}
(b) برهن على أن $e_n \rightharpoonup 0$ بينما $\norm{e_n} = 1$: فالنهايات
الضعيفة قد تفقد المعيار. وبرهن على $\norm x \leq
\liminf\norm{x_{k_j}}$ في (a).
\end{exercise}

\begin{exercise}[$\star\star$]\label{exo:b3:hilbert:8}
(المرافقات) من أجل $T \in \mathcal L(H)$، برهن على وجود $T^* \in \mathcal L(H)$ وحيد يحقق
$\langle Tx, y\rangle = \langle
x, T^*y\rangle$ (بريس)، وعلى $\vertiii{T^*} = \vertiii T$. واحسب مرافق الإزاحة $S$ على
$\ell^2$، وبرهن على $\ker T^* = (\operatorname{im}T)^\perp$ --- واستنتج $\overline{\operatorname{im}T} = (\ker T^*)^\perp$.
\end{exercise}

\begin{exercise}[$\star\star$]\label{exo:b3:hilbert:9}
لتكن $a$ كما في لاكس--ميلغرام وفوق ذلك \emph{متناظرة}. برهن على
أن $u$ تحل $a(u, \cdot) = \varphi$ إذا وفقط إذا كانت $u$ تصغّر $J(v) = \frac12a(v, v) - \varphi(v)$، وعلى أن
الأصغرية مبلوغة عند نقطة واحدة بالضبط. \emph{(بإكمال المربّع:
$J(u + w) - J(u) = \frac12a(w,w) \geq
\frac\alpha2\norm w^2$.)} وتطبيقٌ: استنبط مبرهنة الإسقاط من أجل الفضاءات
الجزئية المغلقة من لاكس--ميلغرام.
\end{exercise}

\begin{exercise}[$\star\star\star$]\label{exo:b3:hilbert:10}
(نظام هار) على $\intcc01$، لتكن $h_{0} = \mathbf 1$، ومن أجل
$n = 2^j + k$ (حيث $j \geq 0$ و $0 \leq k < 2^j$):
\[
h_n = 2^{j/2}\Bigl(\mathbf 1_{[k2^{-j},\,(k +
\frac12)2^{-j})} - \mathbf 1_{[(k+\frac12)2^{-j},\,(k+1)2^{-j})}
\Bigr).
\]
برهن على أن $(h_n)_{n\geq0}$ متعامدة متجانسة في
$L^2(\intcc01)$، وشاملة. \emph{(التعامد: بحوامل منفصلة أو
متداخلة؛ والشمول: تحتوي الفضاءات المولَّدة المنتهية جميعَ الدوال
السلّمية الثنائية، وهي كثيفة --- عبر \cref{thm:b3:lp:density}(1) والتقريب
الثنائي للفترات.)} ونظام هار سلفُ المويجات.
\end{exercise}

\begin{exercise}[$\star\star$]\label{exo:b3:hilbert:11}
(الإسقاطات المتعامدة، مميَّزةً) ليكن $H$ \omterm{def:b3:hilbert:inner}{فضاء هيلبرت} وليكن
$P \in \mathcal L(H)$ يحقق $P^2 = P$ و $P \neq 0$. برهن على تكافؤ:
(أ) أن $P$ هو الإسقاط المتعامد على $\operatorname{im}P$؛
(ب) أن $P = P^*$ (\cref{exo:b3:hilbert:8})؛
(ج) أن $\vertiii P = 1$.
\emph{((ج) $\Rightarrow$ (أ): إذا كانت $x \in
(\ker P)^\perp$ ما تحقق
$Px \neq x$، فلننظر في $x + t(Px - x)$ --- أو مباشرةً: من أجل
$u \in \operatorname{im}P$ و $v \in
\ker P$، انشر $\norm{P(u + tv)}^2 \leq \norm{u + tv}^2$ من أجل كل $t \in \R$ واخلص إلى
$\langle u, v\rangle = 0$.)}
وأبرز إسقاطًا غير متعامد على $\R^2$ واحسب معياره.
\end{exercise}

\begin{exercise}[$\star\star\star$]\label{exo:b3:hilbert:12}
(مبرهنة فون نويمان الإرغودية) ليكن $U \in \mathcal L(H)$ \emph{واحديًّا}
($U^*U = UU^* = I$)، وليكن $F = \ker(U - I)$ فضاء النقاط الصامدة، وليكن $P$ الإسقاط
المتعامد على $F$، ولتكن $A_n = \frac1n\sum_{k=0}^{n-1}U^k$.
(a) برهن على $\ker(U - I) = \ker(U^* - I)$ \emph{(انطلاقًا من $\norm{Ux - x}^2 = 2\norm x^2 - 2\operatorname{Re}\langle
Ux, x\rangle$ ومن الواحدية)}،
واستنتج $\overline{\operatorname{im}(U - I)} = F^\perp$.
(b) برهن على أن $A_nx \to x$ من أجل $x \in F$، وأن
$A_nx \to 0$ من أجل $x \in \operatorname{im}(U - I)$ \emph{(بالتلسكب)}، ثم من أجل
$x \in \overline{\operatorname{im}(U - I)}$ (بالحد المنتظم $\vertiii{A_n} \leq 1$).
(c) اخلص: أن $A_nx \to Px$ من أجل \emph{كل} $x \in H$ --- أي
إن المتوسطات الزمنية تتقارب إلى الإسقاط على الصوامد.
(d) وفصّل ذلك من أجل $H = L^2(\R/\Z)$ و $Uf = f(\cdot +
\alpha)$ حيث $\alpha$
غير ناطق: عيّن $F$ (باستعمال متسلسلات فورييه، \cref{thm:b3:hilbert:fourier})
واستنتج أن $\frac1n\sum_{k<n}f(x + k\alpha) \to \int_0^1f$ في $L^2$: أي التوزّع المتساوي بالمعنى $L^2$
للدورانات غير الناطقة.
\end{exercise}

\section{مسألة: كثيرات الحدود المتعامدة}

\begin{problem}[{مسألة نهاية الأسبوع --- لوجاندر وإرميت
والتربيع الغاوسي}]\label{pb:b3:hilbert:1}
لتكن $I \subseteq \R$ فترةً وليكن $w > 0$ \emph{وزنًا} \omterm{def:b3:topology:continuity}{متصلًا}
على داخل $I$ بحيث يكون $\int_I
\abs t^nw(t)\dd t < \infty$ من أجل كل $n$؛ ونعمل في $H = L^2(I,
w\,\dd\lambda)$
مع $\langle f, g\rangle = \int_I \bar
fg\,w$. ويعطي غرام--شميدت مطبَّقًا على $1, t, t^2, \dots$
\emph{\omterm{pb:b3:hilbert:1}{كثيرات الحدود المتعامدة}} $(p_n)$ الموافقة للوزن $w$
(بالتعيير الموحَّد: $p_n = t^n + \cdots$).\index{كثيرات حدود متعامدة}

\medskip
\noindent\textbf{الجزء الأول --- النظرية العامة.}
\begin{enumerate}
  \item برهن على أن $p_n$ متعامد مع كل كثير حدود درجته $< n$،
        وعلى أن $(p_0, \dots, p_n)$ أساسٌ للفضاء $\R_n[t]$.
  \item (العلاقة التراجعية ذات الحدود الثلاثة) برهن على وجود
        أعداد حقيقية $a_n,
        b_n$ تحقق
        \[
        p_{n+1}(t) = (t - a_n)\,p_n(t) - b_n\,p_{n-1}(t),
        \qquad b_n = \frac{\norm{p_n}^2}{\norm{p_{n-1}}^2} >
        0 .
        \]
        \emph{(انشر $t\,p_n$ في الأساس $(p_k)_{k \leq
        n+1}$ وألغِ المعاملات
        بالتعامد، مستعملًا $\langle tp_n, p_k\rangle = \langle p_n,
        tp_k\rangle$.)}
  \item (الجذور) برهن على أن لكثير الحدود $p_n$ عددًا $n$ من
        الجذور \emph{المتمايزة}، وكلها داخلية في $I$.
        \emph{(لتكن $t_1 < \dots < t_m$ تغيّرات الإشارة الداخلية للمقدار
        $p_n$؛ فإذا كان $m < n$، اختبر $p_n$ مقابل $\prod_{i\leq m}(t - t_i)$
        وناقض التعامد.)}
\end{enumerate}

\medskip
\noindent\textbf{الجزء الثاني --- لوجاندر ($I = \intcc{-1}1$، $w =
1$).}
نعرّف $P_n(t) = \frac{1}{2^nn!}\,\frac{\dd^n}{\dd
t^n}\bigl[(t^2 - 1)^n\bigr]$ (رودريغ).
\begin{enumerate}[resume]
  \item برهن على أن $\deg P_n = n$ بمعامل رئيسي $\frac{(2n)!}{2^n(n!)^2}$، وعلى
        أن $\langle P_n, Q\rangle = 0$ من أجل كل كثير حدود $Q$ درجته $< n$، وذلك
        بالمكاملة بالتجزئة $n$ مرة: فتكون $P_n$ (إلى غاية
        التعيير) \omterm{pb:b3:hilbert:1}{كثيرات الحدود المتعامدة} الموافقة للوزن
        $w = 1$.
  \item احسب $\norm{P_n}_2^2 = \frac{2}{2n+1}$ \emph{(بالمكاملة بالتجزئة $n$ مرة مقابل
        نفسه ثم الإرجاع إلى تكامل بيتا أو واليس،
        \cref{exo:b3:product:8})}.
  \item برهن على أن كثيرات حدود لوجاندر المعيَّرة تشكّل \omterm{def:b3:hilbert:onb}{أساسًا
        هيلبرتيًّا} للفضاء $L^2(\intcc{-1}1)$ \emph{(بفايرشتراس،
        \cref{cor:b3:complete:weierstrass}، مع كثافة $\mathcal C$ في $L^2$)}، وانشر
        $f(t) = \abs t$ حتى الدرجة $2$: أي احسب أفضل تقريب
        تربيعي بالمعنى $L^2$ للمقدار $\abs t$.
\end{enumerate}

\medskip
\noindent\textbf{الجزء الثالث --- إرميت ($I = \R$، $w(t) =
\eu^{-t^2}$).}
نعرّف $H_n(t) =
(-1)^n\eu^{t^2}\frac{\dd^n}{\dd t^n}\eu^{-t^2}$.
\begin{enumerate}[resume]
  \item برهن على أن $H_n$ كثير حدود درجته $n$ بمعامل رئيسي
        $2^n$، وعلى أن $H_{n+1} = 2tH_n -
        H_n'$، وعلى أن $\langle H_m, H_n\rangle_w =
        \delta_{mn}\,2^nn!\sqrt\pi$
        \emph{(بالتجزئة مرة أخرى)}.
  \item برهن على أن عائلة إرميت شاملة في $L^2(\R,
        \eu^{-t^2}\dd t)$، مسلِّمًا
        بنتيجة واحدة من \cref{ch:b3:fouriertransform}: أنه إذا كانت
        $g \in L^1(\R)$ تحقق $\int g(t)\eu^{-\iu\xi t}\dd t = 0$ من أجل كل $\xi$، فإن
        $g = 0$ في كل مكان تقريبًا. \emph{(من أجل $f \perp$ جميع $H_n$،
        أي $\perp$ جميع كثيرات الحدود: برهن على أن
        $z \mapsto \int
        f(t)\eu^{-t^2}\eu^{-\iu zt}\dd t$ معرَّف جيدًا، وانشر الأُسّي في متسلسلة، وبرّر
        التبديل بالهيمنة، واخلص إلى أن تحويل فورييه للمقدار
        $f\eu^{-t^2}$ منعدم.)}
\end{enumerate}

\medskip
\noindent\textbf{الجزء الرابع --- التربيع الغاوسي.} نثبّت $n$،
ولتكن $t_1 < \dots < t_n$ جذور $p_n$ (الجزء الأول)، ونعرّف الأوزان $w_i = \int_I \ell_i(t)\,w(t)\dd t$
حيث $\ell_i$ كثيرات حدود أساس استيفاء لاغرانج عند النقاط
$t_i$.
\begin{enumerate}[resume]
  \item برهن على أن قاعدة التربيع $Q(f) = \sum_iw_if(t_i)$ مضبوطة على كل
        كثيرات الحدود التي درجتها $\leq n - 1$ (بالاستيفاء)،
        بل --- وهذه هي المعجزة --- على كل كثيرات الحدود التي
        درجتها $\leq 2n - 1$: اكتب $P =
        qp_n + r$ واستعمل التعامد
        على خارج القسمة $q$.
  \item برهن على أن الأوزان موجبة \emph{(بتطبيق القاعدة على
        $\ell_i^2$، ودرجته $2n -
        2$)}، واستنتج من مبرهنة
        بوليا (\cref{exo:b3:banach:9}) أن التربيع الغاوسي يتقارب:
        $Q_n(f) \to \int_I fw$ من أجل كل دالة متصلة $f$ على $I$ متراص.
  \item من أجل $n = 2$ و $I = \intcc{-1}1$ و $w = 1$: احسب العقدتين
        $\pm\frac1{\sqrt3}$ والوزنين $1, 1$، وتحقق يدويًّا من الضبط على
        $1, t, t^2, t^3$. وقارن بقاعدة شبه المنحرف على نقطتَي
        التقييم نفسيهما.
\end{enumerate}

\medskip
\noindent\textbf{الجزء الخامس --- تشيبيشيف: كثيرات الحدود
الأفضل تذبذبًا.} والآن $I = \intcc{-1}1$ و $w(t) =
\frac1{\sqrt{1 - t^2}}$.
\begin{enumerate}[resume]
  \item برهن على أن $T_n(\cos\theta) = \cos n\theta$ يعرّف كثير حدود $T_n$ درجته $n$
        (بإثبات $T_{n+1} =
        2t\,T_n - T_{n-1}$ من متطابقة مثلثية)، بمعامل رئيسي
        $2^{n-1}$ من أجل $n \geq 1$؛ وعلى أن التعويض
        $t = \cos\theta$ يعطي
        \[
        \langle T_m, T_n\rangle_w =
        \int_0^\pi\cos m\theta\,\cos n\theta\,\dd\theta
        = 0 \ (m \neq n), \qquad
        \norm{T_0}_w^2 = \pi,\ \ \norm{T_n}_w^2 = \frac\pi2 :
        \]
        فتكون $T_n$ \omterm{pb:b3:hilbert:1}{كثيرات الحدود المتعامدة} الموافقة لهذا
        الوزن، وتكون نشرات تشيبيشيف \emph{هي} متسلسلات فورييه
        الجيبتمامية متنكّرة.
  \item عيّن صراحةً الجذور $n$ $t_k =
        \cos\frac{(2k-1)\pi}{2n}$ والقيم القصوى $n + 1$
        $s_j = \cos\frac{j\pi}n$ للمقدار $T_n$ على $\intcc{-1}1$، حيث $T_n(s_j) = (-1)^j$:
        فالمنحني البياني \emph{يتذبذب بالتساوي} بين $\pm1$.
  \item (الأصغري الأعظمي) برهن على أنه من بين كل كثيرات
        الحدود \emph{الموحَّدة} من الدرجة $n$، يملك كثير
        الحدود $2^{1-n}T_n$ أصغر معيار سوپريموم على
        $\intcc{-1}1$، وهو $2^{1-n}$ --- وهو المصغِّر الوحيد.
        \emph{(فلو كان لكثير حدود موحَّد $P$ أن يحقق $\sup\abs P < 2^{1-n}$،
        لتناوب الفرق $2^{1-n}T_n -
        P$، ودرجته $\leq n-1$، في الإشارة
        عند نقاط التذبذب المتساوي $n+1$.)}
  \item وتطبيقٌ على الاستيفاء: من أجل عقد $t_1 < \dots
        < t_n$ في
        $\intcc{-1}1$، يتضمّن خطأ استيفاء لاغرانج لدالة من
        الصنف $\mathcal C^n$ المقدارَ $\omega(t) = \prod_i(t - t_i)$. برهن على أن
        اختيار جذور تشيبيشيف عقدًا يصغّر $\sup_{\intcc{-1}1}\abs\omega$، وأعطِ الحدّ
        الناتج $\norm{f -
        L_nf}_\infty \leq \frac{\norm{f^{(n)}}_\infty}
        {2^{n-1}\,n!}$ --- وقارنه بالعقد المتساوية التباعد
        (مع ذكر ظاهرة رونغه بوصفها القصة العبرة).
  \item تحقق من $\abs{T_n'(\pm1)} = n^2$ \emph{(باشتقاق $T_n(\cos\theta) = \cos n\theta$ وأخذ النهايات
        $\theta \to 0, \pi$)}: فكثيرات الحدود المحدودة بالعدد $1$ على
        $\intcc{-1}1$ قد تبلغ مشتقتها $n^2$ عند الحافة (وتقول
        متراجحة ماركوف إنها لا تتجاوز ذلك --- والعبارة دون
        برهان). وأين في الفترة يكون حدّ المشتقة $O(n)$ فقط؟
  \item (تربيع تشيبيشيف--غاوس) برهن على أن قاعدة غاوس للوزن
        $w$ عند جذور تشيبيشيف $n$ ذات أوزان \emph{متساوية}
        $w_i = \frac\pi n$ \emph{(بالضبط على $T_0, \dots, T_{n-1}$ مع المجاميع
        المثلثية $\sum_{k=1}^n\cos\bigl(j\tfrac{(2k-1)\pi}{2n}\bigr) =
        0$ من أجل $1 \leq j \leq n - 1$)}: أي أكثر التربيعات
        انتظامًا جميعًا. واكتبها صراحةً من أجل $n = 3$.
\end{enumerate}

\medskip
\noindent\textbf{الجزء السادس --- كريستوفل--داربو، والتداخل،
ومصفوفة ياكوبي.} ونعود إلى وزن عام؛ $h_k = \norm{p_k}^2$ (حيث $p_k$
موحَّدة)، و $b_k = h_k/h_{k-1}$.
\begin{enumerate}[resume]
  \item (أصغر معيار) برهن على أنه من بين كل كثيرات الحدود
        \emph{الموحَّدة} من الدرجة $n$، يكون المتعامد $p_n$
        وحدَه ذا المعيار الأصغر بالمعنى $L^2(w)$ --- وعيّن
        هذا التصغير بوصفه إسقاطًا متعامدًا على $\R_{n-1}[t]$
        (\cref{thm:b3:hilbert:projection} أو الإسقاط المنتهي البُعد من السنة
        الجامعية 2). وخاصية الأصغري الأعظمي في
        السؤال 14 هي العبارة نفسها مع $L^\infty$ بدل $L^2$:
        البطل واحد، والمعياران اثنان.
  \item (كريستوفل--داربو) برهن، بالتراجع على $n$ باستعمال
        العلاقة التراجعية ذات الحدود الثلاثة، على المتطابقة
        \[
        \sum_{k=0}^{n}\frac{p_k(x)\,p_k(y)}{h_k}
        = \frac{p_{n+1}(x)\,p_n(y) -
        p_n(x)\,p_{n+1}(y)}{h_n\,(x - y)}
        \qquad (x \neq y),
        \]
        وعلى صيغتها الالتحامية ($y \to x$):
        $\sum_{k\leq n}\frac{p_k(x)^2}{h_k} =
        \frac{p_{n+1}'(x)p_n(x) - p_n'(x)p_{n+1}(x)}{h_n}$.
  \item استنتج أنه ليس للمقدارين $p_n$ و $p_{n+1}$ أي جذر مشترك،
        وأنه عند كل جذر $x_0$ للمقدار $p_{n+1}$:
        $p_n(x_0)\,p_{n+1}'(x_0) > 0$. واخلص إلى \emph{تداخل} الجذور: أي إن بين كل
        جذرين متتاليين للمقدار $p_{n+1}$ جذرًا واحدًا بالضبط
        للمقدار $p_n$.
  \item (مصفوفة ياكوبي) لتكن $J_n$ المصفوفةَ المتناظرة ثلاثية
        الأقطار من الرتبة $n\times n$ ذات القطر $a_0,
        \dots, a_{n-1}$
        والعناصر خارج القطر $\sqrt{b_1},
        \dots, \sqrt{b_{n-1}}$. برهن بالتراجع على أن
        $\det(tI_n - J_n) = p_n(t)$، ومنه تكون جذور $p_n$ هي القيم الذاتية
        لمصفوفة متناظرة حقيقية --- فيُعاد البرهان في سطر واحد
        على أنها حقيقية، ويُربط (مع التداخل أعلاه) \omterm{pb:b3:hilbert:1}{كثيرات
        الحدود المتعامدة} بالعالم الطيفي في \cref{ch:b3:spectral}.
  \item (تركيب) كوِّن المعجم الخاص بالعائلات الكلاسيكية
        الثلاث (لوجاندر، إرميت، تشيبيشيف): الفترة، والوزن،
        والصيغة المعرِّفة، والعلاقة التراجعية ذات الحدود
        الثلاثة، والمعيار، والموطن الطبيعي لكلٍّ منها (التربيع
        والتقريب على المتراصات؛ والتحليل الغاوسي؛ والأصغري
        الأعظمي وطرق فورييه الجيبتمامية). وجملة واحدة عمّا
        أعطته النظرية العامة (الجزآن الأول والسادس) ولم يكن
        أي حساب مفرد ليعطيه.
\end{enumerate}

\medskip
\noindent\textbf{الجزء السابع --- حدّ الخطأ، والنواة الكامنة
وراء الأوزان.} وهنا يكون $I$ متراصًّا و $f \in \mathcal
C^{2n}(I)$.
\begin{enumerate}[resume]
  \item (صيغة خطأ غاوس) ليكن $Hf$ مستوفي إرميت من الدرجة
        $\leq 2n - 1$ الموافق للدالة $f$ ولمشتقتها $f'$ عند
        العقد $t_1, \dots, t_n$ (وبرهن على وجوده وعلى الخطأ النقطي
        \[
        f(t) - Hf(t) =
        \frac{f^{(2n)}(\xi_t)}{(2n)!}\;p_n(t)^2
        \]
        بحجة الدالة المساعدة المعتادة). واستنتج، بمكاملة هذه
        المتطابقة مقابل $w$ وبالحصر بين القيم القصوى للمقدار
        $f^{(2n)}$، أن
        \[
        \int_I f\,w - Q_n(f)
        = \frac{f^{(2n)}(\xi)}{(2n)!}\,h_n
        \qquad\text{من أجل} \xi \in I,
        \]
        مع $h_n = \norm{p_n}^2$ كما في الجزء السادس: أي إن التربيع الغاوسي
        يخطئ بمشتقة واحدة من الرتبة $2n$، موزونةً بمربّع معيار
        كثير الحدود المتعامد الموحَّد.
  \item (الأوزان قيمٌ لكريستوفل) باستعمال النواة الاستنساخية
        $K_n(x, y) = \sum_{k=0}^{n-1}
        \frac{p_k(x)p_k(y)}{h_k}$ للفضاء $\R_{n-1}[t]$ وضبط $Q_n$ حتى الدرجة
        $2n - 2$، برهن على
        \[
        w_i \;=\;
        \Bigl(\,\sum_{k=0}^{n-1}
        \frac{p_k(t_i)^2}{h_k}\Bigr)^{\!-1} :
        \]
        أي إن كل وزن هو قيمة \emph{دالة كريستوفل} عند عقدته
        --- وهي إيجابية الأوزان (السؤال 10) من جديد، لكن
        بصيغة مضبوطة الآن. وتحقق من أنها تستعيد
        $w_1 = w_2 = 1$ من أجل $n =
        2$ و $I = \intcc{-1}1$ و $w = 1$.
  \item (كل شيء يُتحقَّق منه على تكامل واحد) من أجل وزن
        تشيبيشيف وعقد $n = 3$، احسب طرفَي
        \[
        \int_{-1}^{1}\frac{t^6}{\sqrt{1 - t^2}}\,\dd t
        = \frac{5\pi}{16},
        \qquad
        Q_3(t^6) = \frac{9\pi}{32},
        \]
        فيكون خطأ التربيع $\frac{\pi}{32}$ بالضبط؛ ثم تحقق من
        أن صيغة الخطأ في السؤال 23 تتنبّأ بهذه القيمة بعينها
        (فهنا $f^{(6)} = 6!$ ثابت، و $h_3 = \norm{2^{-2}T_3}_w^2 =
        \frac\pi{32}$): فتتفق النظرية
        والحساب حتى الرقم الأخير.
\end{enumerate}
\end{problem}
